\section{绪论}
\label{sec:01-introduction}

\begin{table}[htbp]
\centering
\begin{tabular}{@{} c c p{8cm} @{}}
\toprule
周次 & 讲次 & 教学内容（要点） \\
\midrule
6 & 1 & 第一章：机器人的起源及发展\newline 第二章：机器人定义及特征 \\
\midrule
\multirow{3}{*}{7} & \multirow{3}{*}{2} & 第三章：机器人坐标系变换、通用旋转变换 \\
& & 第四章：机器人末端位姿描述 \\
& & 第五章：机器人坐标系及D-H法 \\
\midrule
\multirow{2}{*}{8} & \multirow{2}{*}{4} & 第六章：机器人正运动学 \\
& & 第七章：机器人逆运动学 \\
\midrule
\multirow{2}{*}{9} & \multirow{2}{*}{5} & 第八章：微分运动 \\
& & 第九章：雅可比矩阵 \\
\midrule
\multirow{2}{*}{10} & \multirow{2}{*}{6} & 第十章：力雅可比 \\
& & 第十一章：雅可比矩阵的几何性质 \\
\midrule
\multirow{2}{*}{11} & \multirow{2}{*}{7} & 第十二章：机器人刚体动力学 \\
& & 第十三章：机器人连杆动力学（拉格朗日） \\
\midrule
\multirow{2}{*}{12} & \multirow{2}{*}{8} & 第十四章：机器人连杆动力学（牛顿欧拉法） \\
& & 第十五章：机器人连杆动力学（举例） \\
\midrule
\multirow{2}{*}{13} & \multirow{2}{*}{9} & 第十六章：关节空间的轨迹规划 \\
& & 第十七章：笛卡尔空间的轨迹规划 \\
\midrule
\multirow{2}{*}{14} & \multirow{2}{*}{10} & 第十八章：机器人位置控制 \\
& & 第十九章：机器人力位混合控制 \\
\midrule
\multirow{2}{*}{15} & \multirow{2}{*}{11} & 重点复习 \\
& & 重点复习 \\
\bottomrule
\end{tabular}
\caption{机器人学I教学日历（2026年）}
\end{table}

考核形式：

平时作业 (15\%) + 编程作业(20\%) + 实验课 (15\%) + 闭卷考试 (50\%)
